\section{Импульс}
\subsection{Второй закон Ньютона в импульсной форме}
%%%%%%%%%%%%%%%%%%%%%%%%%%%%%%%%%%%%%%%%%%%%%%%%%%%%
%%%%%%%%%%%%%%%%%%%%%%%%%%%%%%%%%%%%%%%%%%%%%%%%%%%%


 \z{Футболист бьет по мячу со средней силой $F=500$ Н. Мяч после удара улетает под углом $45\degree$ к горизонту и приземляется на расстоянии $L=40$ м. Определить время удара по мячу. Сопротивлением воздуха пренебречь. Масса мяча $m=0.5$ кг.}
%
{$t = \frac{m}{F} \sqrt{\frac{L g}{\sin 2\alpha}} = 0.02$ c.}
%%%%%%%%%%%%%%%%%%%%%%%%%%%%%%%%%%%%%%%%%%%%%%%%%%%%
\z{Космический корабль должен, изменив курс, двигаться с прежним по модулю импульсом $p$ под углом $\alpha$ к первоначальному направлению. На какое наименьшее время нужно включить двигатель с силой тяги $F$ и как при этом нужно ориентировать ось двигателя?}
%
{$T = \frac{2p \sin \frac{\alpha}{2}}{F}$, под углом $\beta = \frac{\pi + \alpha}{2}$ к начальной скорости}
%%%%%%%%%%%%%%%%%%%%%%%%%%%%%%%%%%%%%%%%%%%%%%%%%%%%	
 \z[{\begin{wrapfigure}[5]{r}{0.3\textwidth}%
		\vspace{-12pt}%
		\begin{flushright}%
 			\includegraphics[width=0.3\textwidth]{6.1.3}%
		\end{flushright}%
	\end{wrapfigure}%
	}]%
	{На покоящееся тело массы $m_1$ налетает со скоростью $v$ тело массы $m_2$. Сила, возникающая при взаимодействии тел, линейно зависящая от времени, растет от нуля до значения $F_0$ за время $t_0$, а затем равномерно убывает до нуля за то же время $t_0$. Определите скорость тел после взаимодействия, считая, что все движения происходят по одной прямой.}
 %
{$u_1 = \frac{F_0 t_0}{m_1}$, $u_2 = v - \frac{F_0 t_0}{m_2}$}
%%%%%%%%%%%%%%%%%%%%%%%%%%%%%%%%%%%%%%%%%%%%%%%%%%%% 
 \z{Частица, имеющая массу $m = 0.1$ г и начальную скорость $V = 100$ м/с, попадает в область, в которой на неё в течение некоторого времени действует постоянная по модулю и направлению сила $F$. К моменту прекращения действия силы частица приобретает скорость $2V$ в направлении, перпендикулярном первоначальному. Под каким углом к первоначальному направлению движения частицы направлена сила $F$? Какую работу совершила сила $F$ над частицей за время своего действия? Влиянием других сил можно пренебречь.}
%
{$\alpha = \pi - \arctan 2 \approx 115\degree$}
%%%%%%%%%%%%%%%%%%%%%%%%%%%%%%%%%%%%%%%%%%%%%%%%%%%% 
 \z{Ящик с песком массы $M$ лежит на горизонтальной плоскости, коэффициент трения с которой равен $\mu$. Под углом $\alpha$ к вертикали в ящик со скоростью $v$ влетает пуля массы $m$ и почти мгновенно застревает в песке. Через какое время после попадания пули в ящик, он, начав двигаться, остановится? При каком значении $\alpha$ он вообще не сдвинется?}
%
{$t = \frac{m v (\sin \alpha - \mu \cos \alpha)}{\mu (m + M) g}$ при $\tan \alpha > \mu$, при $\tan \alpha \leq \mu$ ящик не сдвинется}
%%%%%%%%%%%%%%%%%%%%%%%%%%%%%%%%%%%%%%%%%%%%%%%%%%%%